\documentclass[oneside,openany,phd]{TMU-Thesis}


% ==================== structure/preamble.tex ====================
% مشخصات پایان‌نامه را در فایلهای faTitle و enTitle وارد نمایید.
% در اینجا برای یکپارچگی، همینجا وارد می‌کنیم

% مشخصات فارسی
\university{تربیت مدرس}
\faculty{مهندسی برق و کامپیوتر}
\department{گروه هوش مصنوعی و رباتیکز}
\subject{مهندسی کامپیوتر}
\field{مهندسی نرم‌افزار}
\credits{---} % در رساله دکتری معمولا به صورت واحدی ذکر نمی‌شود یا فیلد خالی است
\title{یادگیری تقویتی عمیق با شکل‌دهی تصادفی پاداش برای مدیریت عدم قطعیت در سیستم‌های خودوفقی}
\firstsupervisor{دکتر ...}
% \secondsupervisor{استاد راهنمای دوم}
\firstadvisor{دکتر ...}
% \secondadvisor{استاد مشاور دوم}
\name{علیرضا}
\surname{کاویانی‌فر}
\studentID{...}
\thesisdate{بهمن ۱۴۰۴}

% مشخصات انگلیسی
\latinuniversity{Tarbiat Modares University}
\latinfaculty{Faculty of Electrical and Computer Engineering}
\latinsubject{Computer Engineering}
\latinfield{Software Engineering}
\latintitle{Randomized Reward-Shaped Deep Reinforcement Learning for Managing Uncertainty in Self-Adaptive Systems}
\firstlatinsupervisor{Dr. ...}
\firstlatinadvisor{Dr. ...}
\latinname{Alireza}
\latinsurname{Kavianifar}
\latinthesisdate{February 2026}
\latinkeywords{Self-Adaptive Systems, Deep Reinforcement Learning, Randomized Reward Shaping, Uncertainty Management, Formal Verification, UPPAAL-SMC}

% فراخوانی بسته‌ها و دستورات استاندارد
% در این فایل، دستورها و تنظیمات مورد نیاز، آورده شده است.
%-------------------------------------------------------------------------------------------------------------------

% در ورژن جدید زی‌پرشین برای تایپ متن‌های ریاضی، این سه بسته، حتماً باید فراخوانی شود
\usepackage{amsthm,amssymb,amsmath}
% بسته‌ای برای تنطیم حاشیه‌های بالا، پایین، چپ و راست صفحه
\usepackage[top=30mm, bottom=25mm, left=25mm, right=35mm]{geometry}
% بسته‌‌ای برای ظاهر شدن شکل‌ها و تصاویر متن
\usepackage{graphicx}
% بسته‌ای برای رسم کادر
\usepackage{framed} 
% بسته‌‌ای برای چاپ شدن خودکار تعداد صفحات در صفحه «معرفی پایان‌نامه»
\usepackage{lastpage}
% بسته‌ و دستوراتی برای ایجاد لینک‌های رنگی با امکان جهش
% \usepackage[pagebackref=false,colorlinks,linkcolor=black,citecolor=black]{hyperref}
% چنانچه قصد پرینت گرفتن نوشته خود را دارید، خط بالا را غیرفعال و  از دستور زیر استفاده کنید چون در صورت استفاده از دستور زیر‌‌، 
% لینک‌ها به رنگ سیاه ظاهر خواهند شد که برای پرینت گرفتن، مناسب‌تر است
% \usepackage[pagebackref=false]{hyperref}
% بسته‌ لازم برای تنظیم سربرگ‌ها
\usepackage{fancyhdr}
\setlength{\headheight}{16pt}
%
\usepackage{setspace}
\usepackage{algorithm}
\usepackage{algorithmic}
\usepackage{subfigure}
\usepackage{titletoc}


% بسته‌ای برای ظاهر شدن «مراجع» و «نمایه» در فهرست مطالب
\usepackage[nottoc]{tocbibind}
% دستورات مربوط به ایجاد نمایه
\usepackage{makeidx}
\makeindex
% تعریف شیم‌های مورد نیاز برای نسخه‌های جدید بسته bidi در برخی محیط‌ها
% Removal of redundant and potentially broken definitions
\makeatletter
% شیم‌های سازگاری برای نسخه‌های جدید لاتک و بسته‌های bidi/xepersian
\providecommand{\IfClassLoadedT}[2]{\@ifclassloaded{#1}{#2}{}}
\providecommand{\IfClassLoadedF}[2]{\@ifclassloaded{#1}{}{#2}}
\providecommand{\IfClassLoadedTF}[3]{\@ifclassloaded{#1}{#2}{#3}}
\providecommand{\IfPackageLoadedT}[2]{\@ifpackageloaded{#1}{#2}{}}
\providecommand{\IfPackageLoadedF}[2]{\@ifpackageloaded{#1}{}{#2}}
\providecommand{\IfPackageLoadedTF}[3]{\@ifpackageloaded{#1}{#2}{#3}}
\providecommand{\@iffileloaded}[3]{\expandafter\ifx\csname ver@#1\endcsname\relax#3\else#2\fi}
\providecommand{\IfFileLoadedT}[2]{\@iffileloaded{#1}{#2}{}}
\providecommand{\IfFileLoadedF}[2]{\@iffileloaded{#1}{}{#2}}
\providecommand{\IfFileLoadedTF}[3]{\@iffileloaded{#1}{#2}{#3}}
\makeatother

%%%%%%%%%%%%%%%%%%%%%%%%%%
% فراخوانی بسته زی‌پرشین و تعریف قلم فارسی و انگلیسی
\usepackage[pagebackref=false,colorlinks,linkcolor=black,citecolor=black]{hyperref}
\usepackage{notoccite}

\ExplSyntaxOn
\cs_if_exist:NTF \g__hyp_linknestlevel_int { } { \int_new:N \g__hyp_linknestlevel_int }
\cs_if_exist:NTF \l__hyp_annot_GoTo_bool { } { \bool_new:N \l__hyp_annot_GoTo_bool }
\cs_if_exist:NTF \l__hyp_annot_URI_bool { } { \bool_new:N \l__hyp_annot_URI_bool }
\cs_if_exist:NTF \l__hyp_href_url_ismap_bool { } { \bool_new:N \l__hyp_href_url_ismap_bool }
\cs_if_exist:NTF \l__hyp_text_enc_uri_print_tl { } { \tl_new:N \l__hyp_text_enc_uri_print_tl }
\cs_if_exist:NTF \l__hyp_uri_tmpa_tl { } { \tl_new:N \l__hyp_uri_tmpa_tl }

% Defining missing l3pdf/pdfdict commands if not present
\cs_if_exist:NTF \pdfdict_put:nno { } { \cs_new:Npn \pdfdict_put:nno #1#2#3 { } }
\cs_if_exist:NTF \pdfdict_put:nnn { } { \cs_new:Npn \pdfdict_put:nnn #1#2#3 { } }
\cs_if_exist:NTF \pdfdict_use:n { } { \cs_new:Npn \pdfdict_use:n #1 { } }
\cs_if_exist:NTF \pdfannot_dict_put:nne { } { \cs_new:Npn \pdfannot_dict_put:nne #1#2#3 { } }
\cs_if_exist:NTF \pdfannot_link:nen { } { \cs_new:Npn \pdfannot_link:nen #1#2#3 { } }
\cs_if_exist:NTF \__hyp_check_link_nesting:TF { } { \cs_new:Npn \__hyp_check_link_nesting:TF #1#2 { #1 } }
\cs_if_exist:NTF \__hyp_text_pdfstring:eoN { } { \cs_new:Npn \__hyp_text_pdfstring:eoN #1#2#3 { } }
\ExplSyntaxOff

% تعریف دستورات گمشده برای جلوگیری از خطای renewcommand در بسته‌های داخلی
\makeatletter
\providecommand{\foottextfont}{}
\providecommand{\LTRfoottextfont}{}
\providecommand{\RTLfoottextfont}{}
\providecommand{\abstractname}{}
\providecommand{\latinabstract}{}
\providecommand{\AutoPersianDigits}{}
\makeatother

\usepackage{xepersian}
\settextfont[Path=font/,Extension=.ttf,Scale=1,
    BoldFont=XB NiloofarBd,
    ItalicFont=XB NiloofarIt,
    BoldItalicFont=XB NiloofarBdIt
]{XB Niloofar}
\setlatintextfont[Scale=0.9]{Times New Roman}

%%%%%%%%%%%%%%%%%%%%%%%%%%
% چنانچه می‌خواهید اعداد در فرمول‌ها، انگلیسی باشد، خط زیر را غیرفعال کنید
\setdigitfont[Path=font/,Extension=.ttf,Scale=1,
    BoldFont=XB NiloofarBd,
    ItalicFont=XB NiloofarIt,
    BoldItalicFont=XB NiloofarBdIt
]{XB Niloofar}
%%%%%%%%%%%%%%%%%%%%%%%%%%
% تعریف قلم‌های فارسی و انگلیسی اضافی برای استفاده در بعضی از قسمت‌های متن
\defpersianfont\titlefont[Path=font/,Extension=.ttf,Scale=1]{XB Titre}
% \defpersianfont\iranic[Scale=1.1]{XB Zar Oblique}%Italic}%
% \defpersianfont\nastaliq[Scale=1.2]{IranNastaliq}

%%%%%%%%%%%%%%%%%%%%%%%%%%
% دستوری برای حذف کلمه «چکیده»
\renewcommand{\abstractname}{}
% دستوری برای حذف کلمه «abstract»
%\renewcommand{\latinabstract}{}
% دستوری برای تغییر نام کلمه «اثبات» به «برهان»
\renewcommand\proofname{\textbf{برهان}}
% دستوری برای تغییر نام کلمه «کتاب‌نامه» به «مراجع»
\renewcommand{\bibname}{مراجع}
% دستوری برای تعریف واژه‌نامه انگلیسی به فارسی
\newcommand\persiangloss[2]{#1\dotfill\lr{#2}\\}
% دستوری برای تعریف واژه‌نامه فارسی به انگلیسی 
\newcommand\englishgloss[2]{#2\dotfill\lr{#1}\\}
% تعریف دستور جدید «\پ» برای خلاصه‌نویسی جهت نوشتن عبارت «پروژه/پایان‌نامه/رساله»
\newcommand{\پ}{پروژه/پایان‌نامه/رساله }

%\newcommand\BackSlash{\char`\\}

%%%%%%%%%%%%%%%%%%%%%%%%%%
\SepMark{-}

% تعریف و نحوه ظاهر شدن عنوان قضیه‌ها، تعریف‌ها، مثال‌ها و ...
\theoremstyle{definition}
\newtheorem{definition}{تعریف}[section]
\theoremstyle{plain}
\newtheorem{theorem}[definition]{قضیه}
\newtheorem{lemma}[definition]{لم}
\newtheorem{proposition}[definition]{گزاره}
\newtheorem{corollary}[definition]{نتیجه}
\newtheorem{remark}[definition]{ملاحظه}
\theoremstyle{definition}
\newtheorem{example}[definition]{مثال}

%\renewcommand{\theequation}{\thechapter-\arabic{equation}}
%\def\bibname{مراجع}
\numberwithin{algorithm}{chapter}
\def\listalgorithmname{فهرست الگوریتم‌ها}
\def\listfigurename{فهرست تصاویر}
\def\listtablename{فهرست جداول}

%%%%%%%%%%%%%%%%%%%%%%%%%%%%
% دستورهایی برای سفارشی کردن سربرگ صفحات
% \newcommand{\SetHeader}{
% \csname@twosidetrue\endcsname
% \pagestyle{fancy}
% \fancyhf{} 
% \fancyhead[OL,EL]{\thepage}
% \fancyhead[OR]{\small\rightmark}
% \fancyhead[ER]{\small\leftmark}
% \renewcommand{\chaptermark}[1]{%
% \markboth{\thechapter-\ #1}{}}
% }
%%%%%%%%%%%%5
%\def\MATtextbaseline{1.5}
%\renewcommand{\baselinestretch}{\MATtextbaseline}
\doublespacing
%%%%%%%%%%%%%%%%%%%%%%%%%%%%%
% دستوراتی برای اضافه کردن کلمه «فصل» در فهرست مطالب با استفاده از titletoc
\titlecontents{chapter}
  [5em]
  {\addvspace{1em}\bfseries}
  {\contentslabel[\chaptername~\thecontentslabel:]{5em}}
  {\hspace*{-5em}}
  {\titlerule*[0.5pc]{.}\contentspage}

\titlecontents{section}
  [5em]
  {}
  {\contentslabel{2em}}
  {\hspace*{-2em}}
  {\titlerule*[0.5pc]{.}\contentspage}

\titlecontents{subsection}
  [8em]
  {}
  {\contentslabel{3em}}
  {\hspace*{-3em}}
  {\titlerule*[0.5pc]{.}\contentspage}

\titlecontents{subsubsection}
  [12em]
  {}
  {\contentslabel{4em}}
  {\hspace*{-4em}}
  {\titlerule*[0.5pc]{.}\contentspage}

\makeatletter
{\def\thebibliography#1{\chapter*{\refname\@mkboth
   {\uppercase{\refname}}{\uppercase{\refname}}}\list
   {[\arabic{enumi}]}{\settowidth\labelwidth{[#1]}
   \rightmargin\labelwidth
   \advance\rightmargin\labelsep
   \advance\rightmargin\bibindent
   \itemindent -\bibindent
   \listparindent \itemindent
   \parsep \z@
   \usecounter{enumi}}
   \def\newblock{}
   \sloppy
   \sfcode`\.=1000\relax}}
\makeatother


% تنظیمات اختصاصی برای فهرست مطالب (TOC)
\setcounter{secnumdepth}{3}
\setcounter{tocdepth}{3}

\raggedbottom

% تنظیمات نقطه‌چین برای فهرست مطالب
\usepackage{tocloft}
\providecommand{\cftchapdotfill}{\cftdotfill{\cftdotsep}}
\providecommand{\cftsecdotfill}{\cftdotfill{\cftdotsep}}
\providecommand{\cftsubsecdotfill}{\cftdotfill{\cftdotsep}}


% تنظیمات سربرگ برای رفع مشکل نمایش حروف فارسی (جعبه‌های خالی)
\usepackage{fancyhdr}
\pagestyle{fancy}
\fancyhf{}
\fancyhead[L]{\thepage}
\fancyhead[R]{\small\leftmark}
\renewcommand{\chaptermark}[1]{\markboth{\chaptername\ \thechapter.\ #1}{}}
\renewcommand{\sectionmark}[1]{\markright{\thesection.\ #1}}
% حذف حالت ایتالیک/کج که باعث به هم ریختگی فونت فارسی در سربرگ می‌شود
\let\oldheadrule\headrule
\renewcommand{\headrule}{\oldheadrule}
\fancyhead[R]{\small\nouppercase{\leftmark}}



\begin{document}
\AutoPersianDigits
\setstretch{1.5}

\setcounter{chapter}{4}

\chapter{ارزیابی تجربی}
\label{chap:evaluation}
\section{مقدمه}

ارزیابی تجربی نقشی اساسی و تعیین‌کننده در تأیید اثربخشی راهکارهای پیشنهادی در سیستم‌های خودوفقی ایفا می‌کند. به دلیل آنکه معماری نوین \lr{RS-DRL} ادعا می‌کند که می‌تواند بدون نیاز به آموزش‌های پیشرفته و آفلاین، با مدیریت تجربیات ناموفق، بر پدیده بهینه‌های محلی غلبه کند، اعتبارسنجی تجربی این دستاوردها بسیار حیاتی است. در این فصل با بهره‌گیری از سناریوهای متعدد ارزیابی، قابلیت‌های راهکار پیشنهادی را در پلتفرم معتبر شبکه حسگرهای بی‌سیم \lr{DeltaIoT} \cite{gerostathopoulos2019deltaiot} مورد سنجش جامع قرار می‌دهیم. 

این فصل در راستای پاسخگویی به سوالات تحقیق مطرح شده در فصل اول سازاماندهی گردیده و هدف اساسی آن نمایش و اثبات برتری این معماری نسبت به راهکارهای پایه و مرسوم موجود است. ساختار فصل بدین گونه است که در ابتدا اهداف ارزیابی تبیین گشته، سپس تنظیمات، مجموعه داده‌ها و محیط پیاده‌سازی معرفی می‌شوند؛ در گام‌های بعدی ضمن معرفی دقیق معیارهای سنجش، دستاوردها و نتایج به دست آمده در سناریوهای تک‌هدفه و چندهدفه و همچنین تحلیل مقیاس‌پذیری و ارزیابی آماری ارائه خواهد شد.

\section{اهداف ارزیابی و سوالات تحقیق}

انجام این ارزیابی تجربی همسو با پرسش‌های اساسی مطرح شده در فصل پیشین، اهداف مشخصی را دنبال می‌کند که در جدول \ref{tab:eval_rq} نقشه‌برداری شده‌اند:

\begin{table}[h!]
\centering
\caption{هم‌راستایی اهداف ارزیابی با سوالات تحقیق}
\label{tab:eval_rq}
\begin{tabular}{|p{7cm}|p{7cm}|}
\hline
\textbf{سوال تحقیق (\lr{RQ})} & \textbf{رویکرد ارزیابی} \\ \hline
\textbf{سوال اول:} طراحی مکانیزم \lr{DRL} برای خروج از بهینه‌های محلی & مقایسه منحنی‌های همگرایی (شکل‌های \ref{fig:learning_curve_v1} و \ref{fig:learning_curve_v1_1} در این رساله) در برابر روش‌های پایه و ارزیابی عملکرد مجانبی. \\ \hline
\textbf{سوال دوم:} ادغام شکل‌دهی پاداش با بازپخش تجربه & انجام مطالعه ابلیشن (\lr{Ablation Study}) در برابر گونه‌های مختلف از جمله \lr{HER} و \lr{Soft HER}. \\ \hline
\textbf{سوال سوم:} مدیریت فضاهای بزرگ و گسسته تطابق & اجرای تست مقیاس‌پذیری در بسترهای \lr{DeltaIoTv1} خ(۲۱۶ گزینه)، \lr{v1.1} (۱۲۹۶ گزینه) و \lr{v2} (۴۰۹۶ گزینه). \\ \hline
\textbf{سوال چهارم:} مقایسه \lr{RS-DRL} با روش‌های پیشرفته موجود & تقابل با روش‌های استانداردی نظیر \lr{$\epsilon$-greedy DQN}، \lr{DLASER+} و \lr{ML4EAS} (جدول‌های \ref{tab:tts_goals} و \ref{tab:tt_goals} در این رساله). \\ \hline
\end{tabular}
\end{table}

\section{تنظیمات آزمایشگاهی}

پیش از ورود به جزئیات آزمایش‌ها، یادآوری می‌گردد که فرمول‌بندی مسئله یادگیری تقویتی شامل فضای حالت، فضای اقدام و توابع پاداش (که برای اهداف مختلف سیستم از جمله کاهش مصرف انرژی، تاخیر و اتلاف بسته‌ها طراحی شده‌اند)، دقیقاً بر اساس مدل ارائه‌شده در فصل‌های پیشین (به‌ویژه روابط ریاضی مطرح شده در فصل پیشین) بنا شده است. این یکپارچگی تضمین می‌کند که عامل هوشمند با همان قواعد و اهداف نظریِ اثبات‌شده، در محیط شبیه‌سازی آموزش می‌بیند و نیازی به بازتعریف مجدد متغیرها در این فصل نمی‌باشد.

به منظور فراهم نمودن قابلیت بازتولید (\lr{Reproducibility}) آزمایش‌ها، جزئیات پیکربندی مطالعه موردی به دقت بیان می‌گردد. 

\subsection{نمونه‌های مطالعه موردی}
در این پژوهش، از دو بستر آزمایشگاهی متمایز جهت ارزیابی جامع روش پیشنهادی استفاده شده است. سیستم اول، پلتفرم شناخته‌شده \lr{DeltaIoT} در حوزه شبکه‌های حسگر بی‌سیم صنعتی است و سیستم دوم، شبیه‌ساز مأموریت هوایی پهپاد \lr{DARTSim} است که با هدف ارزیابی استقلال از دامنه و مقیاس‌پذیری روش پیشنهادی به کار گرفته شده است. پیکربندی استقراریافته برای شبکه‌های مورد بحث سیستم \lr{DeltaIoT}، مشخصات فضای حالت‌ها و پیچیدگی‌ها بر اساس ویژگی‌های محک \lr{DeltaIoT} \cite{gerostathopoulos2019deltaiot} در جدول \ref{tab:case_study_instances} خلاصه گردیده است:

\begin{table}[h!]
\centering
\caption{نمونه‌های مورد مطالعه سیستم \lr{DeltaIoT} \cite{gerostathopoulos2019deltaiot}}
\label{tab:case_study_instances}
\begin{tabular}{|c|c|c|c|}
\hline
\textbf{نمونه (نسخه)} & \textbf{تعداد سنسورها} & \textbf{اندازه فضای وفقی (گزینه‌ها)} & \textbf{گام‌های آموزش} \\ \hline
\lr{DeltaIoTv1} & 15 & 216 & 10000 \\ \hline
\lr{DeltaIoTv1.1} & 16 & 1296 & 64800 \\ \hline
\lr{DeltaIoTv2} & 37 & 4096 & متناسب با مقیاس شبکه \\ \hline
\end{tabular}
\end{table}

\subsection{مجموعه داده‌ها}
جهت شبیه‌سازی شرایط در زمان اجرا، از ردپاهای محیطی واقعی (\lr{Environmental Traces}) مستخرج از پژوهش جرواستاتوپولوس و همکاران \cite{gerostathopoulos2019deltaiot} مرتبط با سیستم \lr{DeltaIoT} استفاده شده است. این مجموعه داده، الگوهای تداخلات شبکه‌ای متغیر مابین $\text{-40dB}$ تا $\text{+15dB}$ را ثبت کرده و همچنین نوسانات مرتبط با بار ترافیک پیام‌ها (بین ۰ تا ۱۰ پیام به ازای هر سنسور) را در طول زمان شامل می‌شود.

\subsection{پلتفرم پیاده‌سازی}
پلتفرم عملیاتی اجراشده در این تحقیق شامل مجموعه‌ای از پیشرفته‌ترین چارچوب‌های محاسباتی است. برای ایجاد محیط یادگیری تقویتی از کتابخانه استاندارد \lr{Gymnasium} \cite{towers2024gymnasium} بهره گرفته شد و پیاده‌سازی شبکه‌های \lr{DQN} بر روی چارچوب مشهور \lr{Stable-Baselines3} \cite{stable-baselines3} انجام شده است. در دل این پیاده‌سازی، مکانیزم بومی‌شکل‌دهی پاداش تصادفی به صورت مستقیم به درون حافظه بازپخش (\lr{Replay Buffer}) تزریق گشته است. به منظور اعتبارسنجی صوری و ایمنی گزینه‌های تطبیق‌یافته نیز موتور \lr{UPPAAL-SMC} \cite{david2015uppaal} و مؤلفه واسط \lr{ActivFORMS} \cite{weyns2018activforms} نقش رابط‌های کنترلی را ایفا می‌کنند.

\subsection{تنظیمات فراپارامترها \lr{(Hyperparameters)}}
به منظور دستیابی به بهترین عملکرد، فراپارامترهای الگوریتم یادگیری تقویتی از طریق جستجوی شبکه‌ای (\lr{Grid Search}) در طی آزمایش‌ها تنظیم شده‌اند. مهم‌ترین پارامترهای نهایی شامل نرخ یادگیری (\lr{Learning Rate})، اندازه دسته (\lr{Batch Size})، فرکانس به‌روزرسانی شبکه هدف (\lr{Target Network Update Frequency})، و برنامه کاهش نرخ اکتشاف (\lr{$\epsilon$-decay schedule}) می‌باشند. مقادیر دقیق این فراپارامترها پس از ارزیابی‌های اولیه انتخاب شده‌اند تا تعادل مناسبی بین اکتشاف (\lr{Exploration}) و بهره‌برداری (\lr{Exploitation}) برقرار گردد و نتایج ارائه‌شده بر پایه این مقادیرِ بهینه‌شده به دست آمده‌اند.

\section{روش‌های پایه و مقایسه}

جهت ارزیابی عملکرد و تفوق رویکرد \lr{RS-DRL}، آن را با مجموعه وسیعی از متدهای موجود مقایسه کرده‌ایم. مقایسه در سه دسته‌بندی هدف‌گذاری به شرح زیر صورت گرفت:

\begin{itemize}
  \item \textbf{مقاصد کمینه‌سازی (\lr{Minimization Goals}):} برای این منظور مدل در برابر عامل استانداردِ یادگیری پایه، نظیر شبکه‌های کیو عمیق استاندارد \lr{$\epsilon$-greedy DQN} \cite{mnih2015human} و همچنین در برابر ترکیبات \lr{RS-DRL} با استراتژی‌های مشهوری چون بازپخش تجربه با نگاه به گذشته \lr{HER} (\lr{RS-DRL-with-HER}) \cite{andrychowicz2017her} و روش پاداش نرم‌تر \lr{Soft HER} \cite{he2020softher} قرار گرفت. 
  \item \textbf{مقاصد ترکیبی اهداف آستانه‌ای محدود (\lr{TTS Goal}):} برای ارزیابی کیفی، الگوریتم‌های شناخته‌شده‌ای همچون روش کاهش فضای وفقی \lr{DLASER+} \cite{metzger2022realizing} و یک سیستم مرجع (\lr{Reference}) که بر پایه جستجوی کامل و جامع تمامی‌حالت‌ها (\lr{Exhaustive Search}) استوار بود، مورد ارزیابی قرار گرفت.
  \item \textbf{مقاصد آستانه‌ای مطلق (\lr{TT Goal}):} راهکار پیشنهادی در تقابل با مدل‌های یادگیری پیشرفته همچون متدولوژی \lr{ML4EAS} \cite{kachi2023hybrid}، نسخه اولیه \lr{DLASER} \cite{metzger2022realizing}، روش بهبود یافته \lr{DLASER+} \cite{metzger2022realizing} و روش مرجع سنجیده شد.
\end{itemize}

\textbf{نکته کلیدی در خصوص انطباق روش‌های \lr{HER}:} ذکر این نکته حائز اهمیت است که مدل‌های \lr{HER} و \lr{Soft HER} اساساً برای محیط‌های مبتنی بر هدف‌های صریح (\lr{Goal-Conditioned Environments}) طراحی و تنظیم گشته‌اند، در حالی که اهداف کیفی در \lr{DeltaIoT} به صورت کمینه‌سازی پیوسته هستند. با این وجود برای فراهم کردن بستر مقایسه علمی در محیط‌های نرم‌افزاری شبکه همانند \lr{DeltaIoT}، ما این راهکارها را متناسب‌سازی کرده و با در نظر گرفتن آستانه‌های تعریف شده اهداف کیفی (\lr{QoS}) به‌عنوان اهداف جایگزین (\lr{Proxy Goals})، فرآیند تغییر مرجع مجدد (\lr{Relabeling}) را بر اساس آنها شبیه‌سازی کردیم. این رویکردِ استفاده از اهداف جایگزین ضروری بود تا بتوان الگوریتم‌های مبتنی بر \lr{HER} را در سناریوهایی که اساساً بدون هدف صریح هستند، به طور عادلانه ارزیابی نمود.

\section{معیارهای ارزیابی}

در راستای متدولوژی‌های سنجش استاندارد موجود، ارزیابی این رساله بر محور دو گروه دسته‌بندی استوار گشته است:

\subsection{معیارهای عملکرد یادگیری عامل هوشمند}
این معیارها برای بررسی روند یادگیری سیستم استفاده می‌شوند:
\begin{enumerate}
  \item \textbf{عملکرد مجانبی (\lr{Asymptotic Performance}):} این سطح، وضعیت نهایی‌ترین ارزش و کیفیتی است که عامل پس از دریافت آموزش کافی به آن دست می‌یابد. اندازه‌گیری آن به‌طور معمول بر پایه محاسبه معدل پاداش به دست آمده روی تعداد محدودی از گام‌های نهایی چرخه امکان‌پذیر است.
  \item \textbf{زمان تا آستانه (\lr{Time to Threshold}):} بیانگر تعداد گام‌های زمانی است که عامل نیازمند طی کردن است تا به یک مقدار پاداش از پیش تعیین‌شده همگرا گردد. به‌عنوان نمونه آستانه مورد قبول در جداول ارزیابی مقداری معادل $0.9$ در نظر گرفته شده است.
  \item \textbf{عملکرد کل (\lr{Total Performance}):} تابعی برای ارزیابی جامع که از مجذور عملکرد مجانبی در طی زمان کم شده از جمع کل مقادیر پاداش اخذ شده به دست می‌آید. این فرمول به صورت زیر نمایش داده می‌شود:
  $$TP = R_{\infty}\cdot T - \sum_{t=1}^{T} r_t$$
  \textbf{تفسیر این معیار:} به دست آمدن مقادیر مثبت نشانگر آن است که منحنی یادگیری مدل همواره پایین‌تر از سطح آرمانی مجانبی بوده و میل به بهبود و رشد در طول زمان دارد. اما مقادیر منفی حاکی از اوج‌گیری اولیه و سپس افت و همگرایی به سطوح پایین‌تر عملکرد در زمان اجراست.
\end{enumerate}

\subsection{معیارهای عملکرد سیستم}
موارد زیر منعکس‌کننده خروجی نهایی سیستم در ساختار شبکه می‌باشند:
\begin{itemize}
  \item \textbf{تلفات بسته‌های ارسالی (\lr{Packet Loss}):} درصد پیغام‌هایی که با خطا مواجه شده و به مقصد نمی‌رسند.
  \item \textbf{تاخیر ارسال (\lr{Latency}):} مدت زمان صرف‌شده برای انتقال اطلاعات به گره مرکزی.
  \item \textbf{مصرف انرژی (\lr{Energy Consumption}):} مصرف نهایی انرژی در گره‌های سنسورها پس از اعمال وفقی‌سازی.
\end{itemize}

\section{نتایج و تحلیل‌ها}

این بخش به‌عنوان قلب ارزیابی رساله، نتایج استخراج شده از شبیه‌سازی‌ها را تدوین و تفکیک نموده است.

\subsection{منحنی‌های یادگیری و رفتار همگرایی}

روند پیشرفت یادگیری الگوریتم در مقایسه با مدل پایه‌ای به وضوح برتری مکانیزم پیشنهادی را تصدیق می‌کند.
در تحلیل رفتار شبکه با استفاده از \lr{DeltaIoTv1}، مشاهده می‌گردد که رویکرد \lr{RS-DRL} توانسته است ارزش مجانبی بسیار درخشان $0.9905$ را (با خطای انحراف معیار $0.0268$) کسب نموده و در مدت زمان نسبتاً کوتاه ۲۲۶۶ گام (با تحمل خطای ۱۵۰ پیش و پس) به آستانه مطلوب نفوذ کند. در نقطه‌ی مقابل، الگوریتم پایه‌ی حریصانه یعنی \lr{$\epsilon$-greedy DQN} نه تنها نتوانست از حد مجانبی $0.4559$ فراتر رود، بلکه هرگز موفق به فتح آستانه‌ی در نظر گرفته شده نیز نگردید. دلیل بنیادین این توفیق \lr{RS-DRL}، آن است که مکانیزم شکل‌دهی تصادفی، به جای اکتفای تنها بر آموزش پایه، با تخصیص پاداش به تجربیات منحصراً ناموفق اجازه می‌دهد این مدل نه تنها به ثبات بالاتری بهبود یابد، بلکه در بالاترین سطح مجانبی با ثبات مطلوبی مستقر بماند و از بهینه‌های محلی فرار کند.

شکل \ref{fig:learning_curve_v1} منحنی یادگیری الگوریتم‌ها را در فضای \lr{DeltaIoTv1} نشان می‌دهد و برتری همگرایی \lr{RS-DRL} را در برابر روش‌های پایه به وضوح به تصویر می‌کشد.

\begin{figure}[h!]
    \centering
    \includegraphics[width=0.7\textwidth]{figures/multiobjective_deltaiotv1.png}
    \caption{منحنی‌های یادگیری الگوریتم‌ها در محیط \lr{DeltaIoTv1} (پاداش در برابر گام‌ها)}
    \label{fig:learning_curve_v1}
\end{figure}

در فضای بسیار پیچیده‌تر با حضور ۱۲۹۶ گزینه تطبیق‌پذیر در بسته‌ی \lr{DeltaIoTv1.1}، راهکار ما در ۱۹۲۹۶ گام به پاداش عالی $0.8983$ ($\pm 0.0932$) نفوذ پیدا کرد، در حالی که همه مدل‌های رقیب عملکرد کل ضعیف‌تری را ثبت کرده و هرگز به خط آستانه نزدیک نشدند (مطابق با شکل \ref{fig:learning_curve_v1_1}).

\begin{figure}[h!]
    \centering
    \includegraphics[width=0.7\textwidth]{figures/multiobjective_deltaiotv1.1.png}
    \caption{منحنی‌های یادگیری الگوریتم‌ها در محیط \lr{DeltaIoTv1.1}}
    \label{fig:learning_curve_v1_1}
\end{figure}

\subsection{مقایسه عملکرد به ازای هر هدف}

به منظور درک دقیق جزئیات کارکرد عوامل کنترلی روی تک‌تک متغیرهای سیستم، به ارزیابی رویکرد پیشنهاد شده با انواع \lr{HER} پرداختیم. جداول \ref{tab:min_goals_v1} و \ref{tab:min_goals_v1_1} عملکرد مجانبی، زمان تا آستانه و عملکرد کل را برای اهداف مختلف نشان می‌دهند.

\begin{table}[h!]
\centering
\caption{نتایج اهداف کمینه‌سازی در \lr{DeltaIoTv1} (پاداش مجانبی / گام‌ها)}
\label{tab:min_goals_v1}
\resizebox{\textwidth}{!}{
\begin{tabular}{|l|ccc|ccc|ccc|}
\hline
\multirow{2}{*}{\textbf{روش}} & \multicolumn{3}{c|}{\textbf{چندهدفه}} & \multicolumn{3}{c|}{\textbf{مصرف انرژی}} & \multicolumn{3}{c|}{\textbf{تلفات بسته‌ها}} \\ \cline{2-10} 
 & \textbf{مجانبی} & \textbf{زمان تا آستانه} & \textbf{عملکرد کل} & \textbf{مجانبی} & \textbf{زمان تا آستانه} & \textbf{عملکرد کل} & \textbf{مجانبی} & \textbf{زمان تا آستانه} & \textbf{عملکرد کل} \\ \hline
\lr{RS-DRL} & \textbf{$\mathbf{0.99}$} & \textbf{2266} & \textbf{--} & \textbf{$\mathbf{0.97}$} & \textbf{--} & \textbf{--} & \textbf{$\mathbf{0.98}$} & \textbf{--} & \textbf{--} \\
\lr{RS-DRL-HER} & $0.61$ & عدم همگرایی & -- & $0.58$ & عدم همگرایی & -- & -- & -- & -- \\
\lr{RS-DRL-SoftHER} & $0.49$ & عدم همگرایی & -- & $0.39$ & عدم همگرایی & -- & -- & -- & -- \\
\lr{DQN} & $0.45$ & عدم همگرایی & -- & -- & عدم همگرایی & -- & -- & -- & -- \\ \hline
\end{tabular}
}
\end{table}

\begin{table}[h!]
\centering
\caption{نتایج اهداف کمینه‌سازی در \lr{DeltaIoTv1.1}}
\label{tab:min_goals_v1_1}
\resizebox{\textwidth}{!}{
\begin{tabular}{|l|ccc|ccc|ccc|}
\hline
\multirow{2}{*}{\textbf{روش}} & \multicolumn{3}{c|}{\textbf{چندهدفه}} & \multicolumn{3}{c|}{\textbf{مصرف انرژی}} & \multicolumn{3}{c|}{\textbf{تلفات بسته‌ها}} \\ \cline{2-10} 
 & \textbf{مجانبی} & \textbf{زمان تا آستانه} & \textbf{عملکرد کل} & \textbf{مجانبی} & \textbf{زمان تا آستانه} & \textbf{عملکرد کل} & \textbf{مجانبی} & \textbf{زمان تا آستانه} & \textbf{عملکرد کل} \\ \hline
\lr{RS-DRL} & \textbf{$\mathbf{0.89}$} & \textbf{19296} & \textbf{--} & \textbf{$\mathbf{0.93}$} & \textbf{--} & \textbf{--} & \textbf{$\mathbf{0.91}$} & \textbf{--} & \textbf{--} \\
\lr{RS-DRL-HER} & -- & عدم همگرایی & -- & -- & عدم همگرایی & -- & -- & -- & -- \\
\lr{RS-DRL-SoftHER} & -- & عدم همگرایی & -- & -- & عدم همگرایی & -- & -- & -- & -- \\
\lr{DQN} & -- & عدم همگرایی & -- & -- & عدم همگرایی & -- & -- & -- & -- \\ \hline
\end{tabular}
}
\end{table}

تحلیل این اهداف نشان می‌دهد:
\begin{itemize}
  \item \textbf{بهینه‌سازی چندهدفه (\lr{Multi-objective}):} متد ما توانست به پاداشی معادل $0.9905$ (در شبکه اول) و $0.8983$ (در شبکه دوم) دست یابد در حالی که بهترین انواع \lr{HER} منحصراً مقادیری مابین $0.49$ الی $0.61$ را ثبت نمودند. علت این افت فاحش در مدل‌های \lr{HER} را باید در ماهیت وابستگیِ این مدل‌ها به تعاریف و تخصیص اهداف صریح، در مقایسه با پویاییِ اهداف نامشخص اکوسیستمِ \lr{SAS} جستجو نمود.
  \item \textbf{مصرف انرژی:} شاخص‌های ارزیابی، مصرف انرژی \lr{RS-DRL} را با ثبت میانگین $0.9756$ (در \lr{v1}) و $0.9301$ (در \lr{v1.1}) بسیار متمایز از مدل‌های پایه (بازه‌‌ی $0.39$ الی $0.58$) نشان داد.
  \item \textbf{تلفات بسته‌ها (\lr{Packet Loss}):} دستیابی راهکار توسعه‌یافته در نگه‌داری امنیت شبکه‌ها به‌عنوان عاملی در مرز نزدیک به حالت بهینه؛ $0.9850$ در شبکه اولیه و تغییر آن به وضعیت $0.9129$ در شبکه‌ی پیچیده‌تر دوم.
  \item \textbf{تاخیر (\lr{Latency}):} شبکه مجهز به \lr{RS-DRL} موفق‌ترین سرعت رکوردگیری را با معدل $0.9964$ و $0.9021$ منعکس کرد.
\end{itemize}

\subsection{ارزیابی مقایسه‌ای در برابر تمام روش‌های پایه}

ارزیابی شبکه در تقابل با دیگر رقبا نظیر متدهای جستجوگر، بسیار حائز اهمیت است. در جداول \ref{tab:tts_goals} و \ref{tab:tt_goals} نتایج مربوط به اهداف آستانه‌ای خلاصه شده است.

\begin{table}[h!]
\centering
\caption{عملکرد سیستم برای اهداف آستانه‌ای ترکیبی \lr{(TTS Goal)}}
\label{tab:tts_goals}
\begin{tabular}{|l|ccc|}
\hline
\textbf{روش (شبکه \lr{v1})} & \textbf{اتلاف بسته‌ها (\%)} & \textbf{تاخیر (\lr{s})} & \textbf{مصرف انرژی (\lr{C})} \\ \hline
\lr{RS-DRL} & \textbf{$\mathbf{10.578}$} & \textbf{$\mathbf{0.097}$} & $12.882$ \\ 
\lr{DLASER+} & $12.890$ & $0.840$ & $13.020$ \\ 
\lr{Reference} & $14.230$ & $0.100$ & $13.100$ \\ \hline
\end{tabular}
\end{table}

\begin{table}[h!]
\centering
\caption{عملکرد سیستم برای اهداف آستانه‌ای مطلق \lr{(TT Goal)}}
\label{tab:tt_goals}
\begin{tabular}{|l|c|c|c|c|}
\hline
\multirow{2}{*}{\textbf{روش}} & \multicolumn{2}{c|}{\textbf{\lr{DeltaIoTv1}}} & \multicolumn{2}{c|}{\textbf{\lr{DeltaIoTv2}}} \\ \cline{2-5}
 & \textbf{اتلاف بسته‌ها (\%)} & \textbf{تاخیر (\lr{s})} & \textbf{اتلاف بسته‌ها (\%)} & \textbf{تاخیر (\lr{s})} \\ \hline
\lr{RS-DRL} & \textbf{$\mathbf{7.153}$} & -- & \textbf{$\mathbf{9.330}$} & $0.251$ \\
\lr{DLASER+} & $7.980$ & -- & -- & -- \\
\lr{ML4EAS} & -- & -- & $10.560$ & -- \\
\lr{Reference} & -- & -- & -- & \textbf{$\mathbf{0.100}$} \\ \hline
\end{tabular}
\end{table}

در خصوص \textbf{هدف \lr{TTS}}، تحت شبکه اولیه، متد \lr{RS-DRL} از دست رفتگی بسته‌های حاوی پیغام را در مرز $10.578$\% متوقف نمود، در حالی که استفاده از ارزیابی جستجوی مرجع (\lr{Reference}) درصد $14.23$\% و همچنین \lr{DLASER+} بالغ بر $12.89$\% اتلاف داشته‌اند. تاخیر ارسال با نرخ $0.097$ بسیار بهتر از رقیب اصلیِ آن یعنی \lr{DLASER+} (با ثبت $0.840$) کارنامه موفقی از خود برجای گذاشت و متوازن با روش مرجع ($0.100$) بود. میزانِ مصرف انرژی (با مقدار $12.882$) نیز نزدیک به رویکرد مرجع ($13.100$) و \lr{DLASER+} ($13.020$) محاسبه گردید. همچنین در \textbf{هدف \lr{TT}}، راهکار ما در پلتفرم \lr{DeltaIoTv1}، به تلفات بسیار کم شبکه ($7.153$\%) در مقابل بهترین رقیبش یعنی \lr{DLASER+} ($7.980$\%) رسید؛ همچنین در پلتفرم عظیم \lr{v2}، موفق شد مقدار $9.330$\% را برای این فاکتور ثبت نماید در حالی که بهترین سیستم هوش ماشینی فعلی (\lr{ML4EAS}) مقداری معادل $10.560$\% ارائه داد.

تنها تفاوت موجود، تخصیص تاخیر در نسخه دوم بود (که \lr{RS-DRL} تاخیر $0.251$ را در مقایسه با $0.100$ روش مرجع ثبت نمود)؛ که البته یک طراحی مبتنی بر خط‌مشی به جهت اولویت‌بندی انتقال شبکه و کاهش اتلاف به جای کاهش زمان تاخیر قلمداد می‌شود و نه یک شکست الگوریتمی.

\subsection{مطالعه ابلیشن: حساسیت به فاکتور شکل‌دهی $\rho$}

یکی از دستاوردهای علمی‌مهم این متد، تخصیص استراتژیک ضریب مفقود به تجربیات گذشته بود. در الگوریتم، مقدار کنترل‌کننده $\rho$ تعیین می‌کند تا چه تعداد از تجربیات ناکامِ ثبت‌شده در مینی‌بچ (ذخیره‌شده در حافظه‌ی سیستم)، با ضریب عدد یک (پاداش مثبت تصنعی) جایگزین شوند. 

در صورتی که تجربه‌های شکست خورده موجود در حافظه بیشتر از نرخ انتخابی ما باشد، کاملاً به صورت تصادفی تعدادی انتخاب گشته و مجدد شکل‌دهی می‌گردند و در غیر این صورت روی تمامی‌آنها شکل‌دهی پیاده می‌گردد. جهت درک دقیق‌تر، بررسی مقادیر مختلف $\rho$ در قالب شکل \ref{fig:ablation_rho} قابل مشاهده است.

\begin{figure}[h!]
    \centering
    \includegraphics[width=0.7\textwidth]{figures/ablation_rho.pdf}
    \caption{تاثیر مقادیر مختلف ضریب شکل‌دهی $\rho$ بر پاداش مجانبی و عملکرد کل}
    \label{fig:ablation_rho}
\end{figure}

\textbf{نکات تحلیلی پیرامون فاکتور:}
ایجاد یک توازن کلیدی حول فاکتور بسیار مهم است. اگر $\rho$ بیش از حد بالا باشد، عامل با نویز بالا مواجه خواهد شد و عملکرد افت می‌کند. متقابلاً مقادیر بسیار پایین نیز به آن معناست که عامل استراتژی جسورانه‌ای برای خروج از بهینه‌های محلی نخواهد داشت و مشابه شبکه‌های پایه فریز می‌شود. این انتخابِ تصادفی در کنار توزیع منطقی $\rho$، از تناسب‌بیش‌ازحد (\lr{Overfitting}) حافظه نیز ممانعت به عمل می‌آورد.

\subsection{تحلیل مقیاس‌پذیری}

مقایسه نتایج مدل حاکی از پیروزی پایدار آن است. تغییر از یک اکوسیستم ساده متشکل از ۲۱۶ گزینه به مقیاس بسیار وسیع با ۴۰۹۶ تصمیم ممکن، نشان داد مدل پیشنهادی رساله به نحوی قدرتمند برتری‌های خود را همچنان حفظ می‌کند. جدول \ref{tab:scalability} این روند مقیاس‌پذیری را نشان می‌دهد.

\begin{table}[h!]
\centering
\caption{تحلیل مقیاس‌پذیری الگوریتم با بزرگ شدن فضای تطابق}
\label{tab:scalability}
\begin{tabular}{|l|c|c|c|c|}
\hline
\textbf{نسخه \lr{DeltaIoT}} & \textbf{سنسورها} & \textbf{گزینه‌های تطابق} & \textbf{گام‌های آموزش موردنیاز} & \textbf{پاداش مجانبی} \\ \hline
\lr{v1} & $15$ & $216$ & $2266$ & $0.9905$ \\ 
\lr{v1.1} & $16$ & $1296$ & $19296$ & $0.8983$ \\ 
\lr{v2} & $37$ & $4096$ & -- & -- \\ \hline
\end{tabular}
\end{table}

تنها زمان همگراییِ یادگیری (جهش از ۲۲۶۶ به ۱۹۲۹۶ گام در نسخه \lr{v1.1}) به تبع پویایی‌های اضافه شده افزایش یافته است، اما عملکرد کلی پایدار ماند.

\subsection{نتایج و تحلیل‌های تجربی شبیه‌ساز مأموریت پهپاد \lr{DARTSim}}
\label{subsec:experimental_results_dartsim}

جزئیات مربوط به معماری، مؤلفه‌های فیزیکی، فرمول‌بندی مسئله به عنوان فرآیند تصمیم‌گیری مارکوف (\lr{MDP}) و آداپتور انطباق آفلاین این شبیه‌ساز پیشتر در بخش \ref{sec:study_case_dartsim} (صفحه \pageref{sec:study_case_dartsim}) به تفصیل شرح داده شده است. در این بخش، نتایج حاصل از ارزیابی تجربی و عملکرد یادگیری الگوریتم پیشنهادی بر روی این سیستم ارائه می‌گردد.

به منظور ارزیابی استقلال از دامنه روش پیشنهادی رساله (\lr{RS-DRL}) و اثبات کارایی آن در سناریوهای فیزیکی-سایبری کاملاً آفلاین (\lr{Offline RL})، الگوریتم پیشنهادی بر روی داده‌های آفلاین مأموریت پهپاد در شبیه‌ساز \lr{DARTSim} آموزش دیده و عملکرد آن در فرار از نقاط بهینه محلی ارزیابی شد.

\subsubsection{تنظیمات بستر داده آفلاین مأموریت پهپاد}
مجموعه داده جمع‌آوری‌شده شامل \lr{50,000} انتقال حالت با ویژگی‌های \lr{17} بعدی در فضای حالت و فضای عمل گسسته \lr{8}گانه است. داده‌های رفتار پهپاد به شکل تعاملی نامطلوب و با پاداش میانگین به شدت پراکنده ($0.006$) و نرخ موفقیت نهایی مأموریت بسیار ناچیز \lr{2.5\%} ثبت شده‌اند. این توزیع به شدت پراکنده و نامطلوب، یک محیط یادگیری آفلاین بسیار چالش‌برانگیز را پدید آورده است که در آن، عدم کاوش کارآمد فضای عمل به سرعت منجر به سقوط یادگیری مدل‌ها به بهینه‌های محلی بی‌ارزش می‌گردد.

\subsubsection{آموزش و مقایسه عملکرد همگرایی}
الگوریتم پیشنهادی \lr{RS-DRL} با نرخ تصادفی‌سازی شکل‌دهی پاداش $\rho=0.3$ در مواجهه با روش پایه \lr{DQN} (بدون تغییر شکل پاداش) تحت ۳ بذر تصادفی مستقل (شامل سیدهای ۴۲، ۴۳ و ۴۴) به مدت \lr{50,000} گام زمانی گرادیان به صورت کاملاً آفلاین آموزش داده شد. نتایج نهایی به دست آمده در جدول \ref{tab:rsdrl_results} خلاصه شده است.

\begin{table}[ht]
\centering
\caption{مقایسه آماری عملکرد الگوریتم \lr{RS-DRL} و \lr{Baseline DQN} در محیط آفلاین \lr{DARTSim} (میانگین $\pm$ انحراف معیار روی ۳ سید تصادفی مستقل بعد از \lr{50,000} گام آموزشی)}
\label{tab:rsdrl_results}
\resizebox{\textwidth}{!}{%
\begin{tabular}{lccc}
\hline
\textbf{معیار ارزیابی} & \textbf{الگوریتم پیشنهادی \lr{RS-DRL}} & \textbf{روش پایه \lr{DQN}} & \textbf{سطح معناداری آماری} \\
\hline
بیشینه میانگین ارزش کیو (\lr{Max Q-value}) & $5.533 \pm 1.425$ & $1.431 \pm 0.607$ & $p=0.0397^*$ \\
نرخ بهبود ارزش کیو (\%) & \multicolumn{2}{c}{$+286.7\%$} & Cohen's $d=3.75$ \\
\hline
میانگین خطای نهایی آموزش (\lr{TD Loss}) & $0.2709 \pm 0.0122$ & $0.0061 \pm 0.0033$ & --- \\
نرخ شکل‌دهی پاداش واقعی & $0.281$ & غیرقابل‌اعمال & هدف $\rho=0.3$ \\
\hline
\end{tabular}%
}
\end{table}

همان‌طور که در جدول \ref{tab:rsdrl_results} مشهود است، روش پیشنهادی \lr{RS-DRL} موفق به دستیابی به بیشینه ارزش کیو نهایی میانگین \textbf{$5.533$} در مقایسه با \textbf{$1.431$} در روش پایه \lr{DQN} شده است. این دستاورد نمایانگر بهبود خیره‌کننده و چشم‌گیر \textbf{$+286.7\%$} در کیفیت سیاست انطباقی است. ارزیابی با آزمون فرض تی استیودنت دو نمونه‌ای مستقل نشان می‌دهد که این اختلاف از نظر آماری با $p = 0.0397 < 0.05$ کاملاً معنادار است. شاخص اندازه اثر کوهن (\lr{Cohen's d}) برابر \textbf{$3.75$} برآورد شده است که اثبات‌کننده یک تفاوت اندازه اثر فوق‌العاده قوی و کارآمدی مطلق روش در هدایت فرآیند اکتشاف آفلاین است.

نمودارهای همگرایی نهایی، خطای نوسانی آموزش و تحلیل آماری پروسه آموزش در قالب شکل‌های \ref{fig:dartsim_convergence}، \ref{fig:dartsim_seeds_breakdown} و \ref{fig:dartsim_statistical_analysis} ترسیم و ارائه شده‌اند.

\begin{figure}[htbp]
  \centering
  \includegraphics[width=0.85\textwidth]{figures/figure1_main_convergence.png}
  \caption{نمودار همگرایی مأموریت پهپادی \lr{DARTSim} شامل چهار پانل: (الف) همگرایی ارزش کیو، (ب) میانگین خطای \lr{TD}، (ج) نرخ شکل‌دهی پاداش واقعی، و (د) مقایسه گام‌به‌گام ارزش کیو.}
  \label{fig:dartsim_convergence}
\end{figure}

\begin{figure}[htbp]
  \centering
  \includegraphics[width=0.85\textwidth]{figures/figure2_per_seed_breakdown.png}
  \caption{تحلیل تفکیکی رفتار الگوریتم \lr{RS-DRL} در مقایسه با \lr{DQN} پایه به ازای سیدهای تصادفی مستقل ۴۲، ۴۳ و ۴۴ در محیط شبیه‌ساز پهپادی \lr{DARTSim}.}
  \label{fig:dartsim_seeds_breakdown}
\end{figure}

\begin{figure}[htbp]
  \centering
  \includegraphics[width=0.85\textwidth]{figures/figure3_statistical_analysis.png}
  \caption{تحلیل آماری و معناداری اندازه اثر کوهن \lr{(Cohen's d)} و فواصل اطمینان ارزش کیو برای روش پیشنهادی در برابر روش پایه در مأموریت‌های \lr{DARTSim}.}
  \label{fig:dartsim_statistical_analysis}
\end{figure}

\subsubsection{تحلیل پدیده نوسان خطا و خوش‌بینی در یادگیری}
در جدول فوق، میانگین خطای \lr{TD Loss} نهایی برای الگوریتم پیشنهادی \lr{RS-DRL} مقدار $0.2709$ در مقایسه با خطای ناچیز $0.0061$ برای \lr{DQN} پایه است. این پدیده کاملاً همسو با تئوری کارکرد شکل‌دهی تصادفی پاداش است. روش پیشنهادی با تبدیل پاداش شکست‌های گذشته به مقدار خوش‌بینانه $1.0$ با نرخ شکل‌دهی پاداش واقعی $0.281$، باعث ایجاد اختلاف مستمر در تخمین‌های شبکه هدف و شبکه جاری می‌شود. این نوسان خطا به عنوان یک رگولارایزر هوشمند عمل کرده، وزن‌های شبکه را از تثبیت زودهنگام در بهینه‌های محلی مصون داشته و همگرایی پایدار به سمت سیاست بهینه سراسری را در فضاهای گسسته و بزرگ تلاطم محیطی تضمین می‌کند.

\section{اعتبارسنجی آماری}

برای استواری یافته‌های این رساله، نتایج حاصله توسط آزمون‌های دقیق آماری اعتبارسنجی شده‌اند. 
روش‌شناسی توسعه‌یافته در ارزیابی شامل حداقل ۳۰ دور اجرای کاملاً مستقل با بذرهای رندوم، انجام آزمون تی تک‌نمونه‌ای (\lr{One-sample t-test}) \cite{field2017discovering, casella2002statistical} با ضریب خطای آلفای برابر $0.05$، و تخصیص فواصل اطمینان ۹۵٪ به روش بوت‌استرپ (\lr{Bootstrap}) \cite{efron1993bootstrap} است.
فرمول ریاضی محاسبه آماری به این شکل تنظیم گردید:
$$t = \frac{\bar{x} - \mu_0}{\frac{s}{\sqrt{n}}}$$

\textbf{خلاصه اعتبارسنجی‌ها (بر اساس مقادیر $p$):}
از بین مجموعاً ۲۸ مقایسه‌ی متمایز و شاخص، تعداد ۲۵ مقایسه کیفیتِ الگوریتم توسعه‌یافته در یک شرایط آماری فوق‌العاده قرار داشتند ($89.3$\% تست‌ها معنادار بودند). جدول \ref{tab:p_values} گزیده‌ای از نتایج آزمون آماری و مقادیر احتمال را نشان می‌دهد.

\begin{table}[h!]
\centering
\caption{مقادیر احتمال \lr{(p-values)} برای مقایسه‌های معنادار و غیرمعنادار}
\label{tab:p_values}
\begin{tabular}{|l|l|c|c|}
\hline
\textbf{معیار (هدف)} & \textbf{رقیب پایه} & \textbf{مقدار $p$ (\lr{v1})} & \textbf{وضعیت معناداری ($p < 0.05$)} \\ \hline
تلفات بسته (\lr{TTS}) & \lr{DLASER+} & $< 0.001$ & معنادار (\checkmark) \\
تلفات بسته (\lr{TTS}) & \lr{Reference} & $< 0.001$ & معنادار (\checkmark) \\
مصرف انرژی (\lr{TTS}) & \lr{Reference} & $0.326$ & غیرمعنادار ($\times$) \\
تاخیر (\lr{TTS}) & \lr{Reference} & $0.134$ & غیرمعنادار ($\times$) \\ \hline
\end{tabular}
\end{table}

شگفتی زمانی کامل گشت که در پارامتر اتلاف بسته‌های شبکه (\lr{Packet Loss}) تمام ۱۰۰٪ از تست‌ها ارزش بهبود معنادار ارائه دادند؛ به‌طور نمونه، احتمال وقوع تصادفی برتری آن برای متد \lr{TTS} شبکه اول در مقابل مدل‌های \lr{DLASER+} و رویکرد مرجع کمتر از $0.001$ برآورد گردید (\lr{$p < 0.001$}). مواردی جزئی نیز که نتایج قابل توجهی نداشتند (مانند آزمون تاخیر که مقدار \lr{$p=0.134$} در شبکه اولیه دریافت نمود یا مصرف انرژی که در حالت \lr{TTS} مقداری معادل \lr{$p=0.326$} ثبت کرد) عمدتاً به دلیل عملکرد متعادل با روش مرجع بوده و هیچگونه شواهدی مبنی بر تضعیف متدِ موردِ بررسی به دست نمی‌دهد.

\section{خلاصه فصل}

این فصل به نحو جامعی توانست ادعاهای نوآوری رساله پیرامون مکانیزم شکل‌دهی یادگیری را تایید نماید. با مرور نتایج مشخص گردید: ۱) متد ابداعی \lr{RS-DRL} با پایداری ویژه‌ای توانسته تمامی‌الگوریتم‌های پرآوازه و قدرتمند فعلی پیشین را در ابعاد بهبود عملکرد یادگیری کیفیت و کنترل سیستم شکست دهد. ۲) این برتری در حدود $89.3$\% مواقع توسط تست‌های اعتبارسنجی آماری تأیید گردیده‌اند (۲۵ مورد از ۲۸ مقایسه کلیدی). ۳) این متد ثابت نمود تغییر مقیاس گزینه‌های وفقی از ۲۱۶ به ۴۰۹۶ گزینه نمی‌تواند قابلیت‌های آن را مختل نماید. ۴) مکانیزم شکل‌دهی مجدد توانسته در مقایسه دقیق و عمیق، از وقوع توقف سیستم‌ها در تله‌های استثناهای محلی جلوگیری کند و نهایتاً ۵) برخلاف متدهای رایج هم‌نوع پیشرفته خود نظیر مدل‌های \lr{HER}، برای پویاییِ پایدار اساساً نیازمند تخصیص دانش سیستم‌ها با دامنه نامشخص (\lr{Domain-specific knowledge}) و یا اهداف صریح از پیش مقدر نبوده است.
\newpage


% ==================== chapters/chapter6.tex ====================


\bibliographystyle{ieeetr-fa}
\bibliography{References,MyReferences}
\end{document}
